\documentclass{article}
\usepackage[T2A]{fontenc}
\usepackage[utf8]{inputenc}
\usepackage[english,russian]{babel}
\usepackage{amsmath, amsfonts, amsthm}
\usepackage{fancyvrb}
\usepackage{indentfirst}
\usepackage[
    left=2cm, right=2cm, top=2cm, bottom=2cm,
    headsep=0.2cm, footskip=0.6cm, bindingoffset=0cm
]{geometry}

\begin{document}

\section*{Вариант 5}

Производя подстановку для одновременного применения обоих сдвигов, получим:
\begin{equation}
\begin{cases}
A = 3\delta - b_s = -b - 3a\Delta + 3\delta - 6\Delta^2, \\[4pt]
B = 3\delta^2 - 2b_s\delta + a_s c_s - 4d_s = \\[4pt]
\quad = 3\delta^2 - 2(b + 3a\Delta + 6\Delta^2)\delta + (a + 4\Delta)(c + 2b\Delta + 3a\Delta^2 + 4\Delta^3) - \\[2pt]
\quad \quad - 4(d + c\Delta + b\Delta^2 + a\Delta^3 + \Delta^4), \\[8pt]
C = \delta^3 - b_s\delta^2 + (a_s c_s - 4d_s)\delta - a_s^2 d_s - c_s^2 + 4b_s d_s = \\[4pt]
\quad = \delta^3 - (b + 3a\Delta + 6\Delta^2)\delta^2 + \Bigl((a + 4\Delta)(c + 2b\Delta + 3a\Delta^2 + 4\Delta^3) - \\[2pt]
\quad \quad - 4(d + c\Delta + b\Delta^2 + a\Delta^3 + \Delta^4)\Bigr)\delta - (a + 4\Delta)^2(d + c\Delta + b\Delta^2 + a\Delta^3 + \Delta^4) - \\[2pt]
\quad \quad - (c + 2b\Delta + 3a\Delta^2 + 4\Delta^3)^2 + 4(b + 3a\Delta + 6\Delta^2)(d + c\Delta + b\Delta^2 + a\Delta^3 + \Delta^4).
\end{cases}
\end{equation}
Расчёты показывают, что когда
\begin{equation}
\delta = \alpha a_s^2 + \beta b_s, \quad \text{где} \quad 8\alpha + 3\beta = 1,
\end{equation}
коэффициенты резольвенты не меняются при любом сдвиге $x$.

\end{document}